\chapter{Single Fluid ideal MHD equations}

This theory follows mostly the EPFL Intro to Plasma Phyiscs course, Freidberg's book \emph{Ideal MHD} 
(\cite{Freidberg2014}) and Chen's book \emph{Introduction to Plasma Physics and Controlled Fusion} (\cite{Chen2016}).

\section{The Maxwell equations}

\begin{align}
    \text{Gauss's law: } \nabla \cdot \vec{E} &= \frac{\rho}{\epsilon_0} \label{eq:gauss_e} \\
    \text{Gauss's law for magnetism: } \nabla \cdot \vec{B} &= 0 \label{eq:gauss_b} \\
    \text{Faraday's law of induction: } \nabla \wedge \vec{E} &= -\frac{\partial \vec{B}}{\partial t} \label{eq:faraday} \\
    \text{Ampère–Maxwell law: }  \nabla \wedge \vec{B} &= \mu_0 \vec{j} + \mu_0 \epsilon_0 \frac{\partial \vec{E}}{\partial t} \label{eq:ampere}
\end{align}

\section{The Boltzmann equation}
The Boltzmann equation is
\begin{equation}
    \frac{\partial f_s}{\partial t} + \vec{v} \nabla_{\vec{r}} f_s + \frac{q_s}{m_s}(\vec{E}^{\lr} \vec{v}\wedge \vec{B}^{\lr} ) \nabla_{\vec{v}} f_s - \left(\frac{\partial f}{\partial t}\right)_{\operatorname{coll}} = 0
\end{equation}

Here the short range interactions are accounted for in the collisional term, and the long range (lr) interactions from collective effeects are accounted for in the electric and Lorentz term.
The crititical distance hereby is the Debye length of electrons $\lambda_{D} = \sqrt{\frac{\epsilon_0 T_e}{e^2 n_0}}$.

\subsection{Gravity irrelevant}

Gravity on single ion: $F_G = m g \sim \SI{1e-26}{\kilo\gram} \SI{10}{\meter\second^{-2}} \sim \SI{1e-25}{\newton}$


From Poisson eq ($\nabla^2 \phi = \frac{\rho}{\epsilon_0}$) and length scale estimate for long range forces ($\nabla \sim \frac{1}{\lambda_D}$), we can estimate the 
long-range (beyond $\lambda_D$) Coulomb force: $F_C = e E = e \nabla \phi \sim \lambda_D \frac{\rho}{\epsilon_0} = \frac{e^2 n_e \lambda_D}{\epsilon_0} \sim \SI{3e-27}{m^2 \newton} n_e \lambda_D$

Even a very weak estimate ($\lambda_D \sim 10^{-6}$, $n_e \sim 10^{14}$) , compared to typical plasma values (see \cite{Blandford2013}), this shows that gravity can be neglected.

\section{Species continuity equation and momentum equation}
\subsection{Species variables}
\begin{itemize}
    \item Number density: $n_s(\vec{r}, t) = \int d\vec{v} \, f_s(\vec{r},\vec{v},t)$
\end{itemize}
In general, we note the per particle average in velocity space of a function $g(\vec{v})$ at a given $(\vec{r},t)$ as
\begin{equation}
    \langle g(\vec{r},t)\rangle_{\vec{v}} = \frac{1}{n_s(\vec{r}, t)} \int d\vec{v} \, g(\vec{v}) f_s(\vec{r},\vec{v},t)
\end{equation}

Thus we have 

\begin{itemize}
    \item Average fluid velocity $\vec{u}_s(\vec{r},t)$: $g(\vec{v})=\vec{v} \rightarrow \vec{u}_s(\vec{r},t) =\frac{1}{n_s(\vec{r}, t)} \int d\vec{v} \, \vec{v} f_s(\vec{r},\vec{v},t)$
    \item Average kinetic energy density $\vec{w}_s(\vec{r},t)$: $g(\vec{v})=\frac{1}{2}m_s\vec{v}^2 \rightarrow \vec{u}_s(\vec{r},t) =\frac{1}{n_s(\vec{r}, t)} \int d\vec{v} \, \frac{1}{2}m_s\vec{v}^2 f_s(\vec{r},\vec{v},t)$
\end{itemize}

Furthermore, we define the pressure tensor $\overleftrightarrow{P}$ as (note that we do not normalize over number density)
\[ \overleftrightarrow{P}_s \coloneq \int d\vec{v} \, (\vec{v}-\vec{u}_s) \otimes (\vec{v}-\vec{u}_s) f_s(\vec{r},\vec{v},t)
\]

We can evaluate moments of the Boltzmann equation


\[
\int d\vec{v} \, g_s(\vec{v}) \left[\text{Boltzmann}\right] = 0
\]
\subsection{Species continuity equation}
$\to$ continuity eq.: choose $g_s(\vec{v}) = 1$

\begin{align*}
&\Rightarrow \int d\vec{v} \, \left[\underbrace{\frac{\partial f_s}{\partial t}}_{(1)} + \underbrace{\vec{v} \cdot \frac{\partial f_s}{\partial \vec{r}}}_{(2)} + \underbrace{\frac{q_s}{m_s} (\vec{E}^{lr} + \vec{v} \wedge \vec{B}^{lr}) \cdot \frac{\partial f_s}{\partial \vec{v}}}_{(3)} - \underbrace{\left(\frac{\partial f}{\partial t}\right)_{coll}}_{(4)}\right] = 0
\end{align*}

\begin{align*}
(1): \quad & \int d\vec{v} \,  \frac{\partial f_s}{\partial t} = \frac{\partial}{\partial t} \int f_s = \frac{\partial}{\partial t} n_s \\
(2): \quad & \int d\vec{v} \,  \vec{v} \cdot \frac{\partial f_s}{\partial \vec{r}} = \frac{\partial}{\partial \vec{r}} \cdot \int d\vec{v} \, \vec{v} f_s = \frac{\partial}{\partial \vec{r}} \cdot (n_s \vec{u}_s) \qquad \text{$\vec{v}$ indep. of $\vec{r}$} \\
& \qquad \qquad \qquad \qquad \qquad \qquad \qquad \text{$\vec{v}$ indep. of $\vec{r}$, $\vec{u} = \langle \vec{v} \rangle$} \\
(3): \quad & \int d\vec{v} \,  \frac{q_s}{m_s} (\vec{E}^{lr} + \vec{v} \wedge \vec{B}^{lr}) \cdot \frac{\partial f_s}{\partial \vec{v}} = \frac{q_s}{m_s} \int \frac{\partial}{\partial \vec{v}} \cdot (\vec{E} + \vec{v} \wedge \vec{B}) \\
& = \frac{q_s}{m_s} \int_{S_{\vec{v} \to \infty}} (\vec{E} + \vec{v} \wedge \vec{B}) \, f_s \, d\vec{S}_{\vec{v}} = 0 \qquad\text{Ass. } f_s\overset{v\to\infty}{\rightarrow}0\\
(4): \quad & \int d\vec{v} \,  \left(\frac{\partial f}{\partial t}\right)_{coll} = 0 \quad \text{bec. we ass. collisions don't create or destroy ptcls.}
\end{align*}

We thus arrive at the species fluid continuity equation
\begin{equation}
    \frac{\partial}{\partial t} n_s + \frac{\partial}{\partial \vec{r}} \cdot (n_s \vec{u}_s) = 0
\end{equation}



\subsection{Species momentum equation}
$\to$ momentum eq.: choose $g_s(\vec{v}) = m_s \vec{v}$

See derivation in \autoref{sec:Species Momentum Equation}

\begin{equation}\label{eq: species momentum eq}
 m_s n_s \frac{d \vec{u}_s}{d t} = q_s n_s \left(\vec{E} + \vec{u}_s \wedge \vec{B}\right) - \frac{\partial}{\partial \vec{r}} \cdot \overleftrightarrow{P_s} + \vec{R}_s
\end{equation}

where $\frac{d}{dt}$ is the convective derivative, $\frac{d}{dt} = \frac{\partial}{\partial t} + \vec{u}_s \cdot \frac{\partial}{\partial \vec{r}}$

$\vec{R}_s$ is a viscous force due to collisions, $\vec{R} = \int d\vec{v} \, m_s (\vec{v} - \vec{u}_s) \left(\frac{\partial f}{\partial t}\right)_{coll} $.

In the relaxation time approcimation we have $\vec{R}_e \approx -m_e n_e \nu_{ei} (\vec{u}_e - \vec{u}_i)$,
 where $\nu_{ei}$ is the electron-ion collision frequency.

\subsection{Species energy equation}
$\to$ energy eq.: choose $g_s(\vec{v}) = m_s \frac{\vec{v}^2}{2}$

{\color{magenta} Show derivation in appendix}


\begin{align*}
\Rightarrow \quad \frac{3}{2} n_s \frac{d T_s}{d t} + \overleftrightarrow{P_s} : \frac{\partial}{\partial \vec{r}} \vec{u}_s + \frac{\partial}{\partial \vec{r}} \cdot \vec{q}_s &= Q_s
\end{align*}

\begin{itemize}
\item double contraction ":" : $\overleftrightarrow{P_s} : \frac{\partial}{\partial \vec{r}} \vec{u}_s \coloneq \sum_{i,j} P_{(s),\, ij} \frac{\partial u_{(s),\, i}}{\partial r_j}$
\item $T_s := \frac{1}{3n_s} \int d\vec{v} \, m_s \, (\vec{v} - \vec{u}_s)^2 f_s$ \quad, \quad the kinetic energy due to spreading around fluid velocity

\item $\vec{q}_s = \frac{1}{2} m_s \int d\vec{v} \, (\vec{v} - \vec{u}_s)^2 (\vec{v} - \vec{u}_s) f_s$ \quad, \quad the heat flux vector, so the kinetic energy that is transported with velocity of the ptcls

\item $Q_s = \int d\vec{v} \, \frac{1}{2} m_s \, (\vec{v} - \vec{u}_s)^2 \left(\frac{\partial f}{\partial t}\right)_{coll}$ \quad, \quad heat generated due to viscous forces
\end{itemize}

$\to$ Can generally evolve each species, often we care about ions \& electrons $\to$ call it two-fluid model



\section{Single fluid model}
\subsection{Single fluid variables}
\begin{itemize}
    \item Mass density: $\varrho(\vec{r}, t) = \sum_s n_s m_s$
    \item Center of Mass Velocity: $\vec{V}(\vec{r}, t) = \frac{\sum m_s n_s \vec{u}_s}{\varrho}$
    \item Total electric charge density: $\rho = \sum_s n_s q_s$
    \item Total electric current density: $\vec{J} = \sum_s q_s n_s \vec{u}_s$
    \item Pressure tensor:
    \begin{itemize}
        \item Species pressure tensor in center of mass frame: \[\overleftrightarrow{P}_s^{cm} = \int d\vec{v} \,m_s\, (\vec{v}-\vec{V}) \otimes (\vec{v}-\vec{V}) f_s(\vec{r},\vec{v},t) \]
        \item Singe fluid pressure tensor in center of mass frame: $\overleftrightarrow{P}^{cm} = \sum_s \overleftrightarrow{P}_s^{cm}$
    \end{itemize} 
\end{itemize}

\subsection{Single fluid equations}

In the following we want to combine the species equations to a model describing the behaviour of the combined fluid consisting of the different species.

\subsubsection{Mass continuity equation}
Multiplying the species continuity equation by their respective mass $m_s$, which is time independent, and adding the the equations for all present species, we obtain the \emph{single fluid mass continuity equation}

\begin{equation}
    \frac{\partial \varrho}{\partial t} + \nabla \cdot (\varrho \vec{V}) = 0
\end{equation}


\subsubsection{Charge continuity equation}
Similarly, multiplying the species continuity equation by their respective charges $q_s$, which is time independent, and adding the the equations for all present species, we obtain the \emph{single fluid charge continuity equation}

\begin{equation}
    \frac{\partial \rho}{\partial t} + \nabla \cdot  \vec{J} = 0
\end{equation} 

\subsubsection{Momentum equation}

Summing the species momentum equations \eqref{eq: species momentum eq}, we can obtain the \emph{single fluid momentum equation}.

\begin{equation}\label{eq: Single fluid momentum eq.}
    \varrho \frac{d\vec{V}}{dt} = \rho \vec{E} + \vec{J} \wedge \vec{B} - \partial_{\vec{r}} \cdot \overleftrightarrow{P}^{cm}
\end{equation}

Details of the derivation can be found in \ref{sec:Derivation of the Single Fluid Momentum Equation}.



\subsubsection{Ohm's law}

{\color{magenta} Clean derivation}

from momentum equations \quad $m_s n_s \frac{d\vec{u}_s}{dt} = q_s n_s \left(\vec{E} + \vec{u}_s \wedge \vec{B}\right) - \partial_{\vec{r}} \cdot \overleftrightarrow{P_s} + \vec{R}_s$

$\to$ Ass. we have electrons and one species of ions

$\to$ multiply electron eq. by $\frac{1}{q_e}$ and ion eq. by $\frac{1}{q_i} = \frac{m_e}{m_i n_e}$

$\to$ assume neutrality $(n_e \simeq n_i \simeq n)$, $m_e \ll m_i$, $\vec{J} \ll e n_e \vec{u}_e$, $e n_i \vec{u}_i$

$\to$ Sum of their respective eq's is approx.:
\[\frac{m_e}{n e^2} \left(\frac{\partial \vec{J}}{\partial t} + \frac{\partial}{\partial \vec{r}} \cdot (\vec{V}\vec{J} + \vec{J}\vec{V})\right) = \vec{E} + \vec{V} \wedge \vec{B} - \frac{1}{en} \frac{\partial}{\partial \vec{r}} \cdot \overleftrightarrow{P}_e^{cm} - \eta \vec{J}\]

with the \emph{plasma resisitivity} $\eta = \frac{m_e \nu_{ei}}{e^2 n} $.
\subsubsection{Closure}

In order to have a closed system of equations, we need additional assumptions and thus restrictions. This is called a \emph{closure}.
Many types of closures exist and most are quite involved.

For an isotropic pressure tensor and adiabatic processes we have the following closure:

\[
\frac{\partial}{\partial \vec{r}} \cdot \overleftrightarrow{P}^{CM} = \frac{\partial}{\partial \vec{r}} P \quad, \quad \frac{d}{dt} (P \varrho^{-\gamma}) = 0 \quad \text{with} \quad \gamma = \frac{c_p}{c_v}= \frac{5}{3}
\]

\begin{comment}
    Three closures

    \[
    \frac{\partial}{\partial \vec{r}} \cdot \overleftrightarrow{P}^{CM} = \frac{\partial}{\partial \vec{r}} P \quad, \quad \frac{d}{dt} (P \varrho^{-\gamma}) = 0 \quad \text{with} \quad \begin{cases} \gamma = 1 & \text{isothermal} \\ \gamma = \frac{c_p}{c_v} & \text{adiabatic} \\ \gamma = +\infty & \text{incompressible} \end{cases}
    \]
\end{comment}

We will justify these assumptions in the next section.


\subsubsection{Single fluid Maxwell equations}
Maxwell equations
\begin{align*}
\nabla \cdot \vec{E} &= \frac{\rho}{\epsilon_0} & \nabla \wedge \vec{E} &= -\frac{\partial \vec{B}}{\partial t} \\
\nabla \cdot \vec{B} &= 0 & \nabla \wedge \vec{B} &= \mu_0 \left(\vec{J} + \epsilon_0 \frac{\partial \vec{E}}{\partial t}\right)
\end{align*}








\section{Magnetohydrodynamic model}

\subsection{Assumptions}

By comparing constants as well as time and length scales typical in fusion plasmas or provided specifically for polomacs by \cite{Elio2024}, we can make the following assumptions.

\begin{itemize}
    \item Characteristic length: $L \sim \SI{1}{\meter}$
    \item We assume approximate thermal equilibrium: $T_e \approx T_i \approx \SI{100}{\electronvolt}$
\end{itemize}

In the following we specifically apply  \cite{Freidberg2014}, chapter 2.

\begin{enumerate}

\item Electron inertia negligible, $m_e \to 0$

Electrons respond nearly instantaneously to field changes on MHD timescales. 
The electron cyclotron period $\tau_{ce} = 2\pi/\omega_{ce}$ is much shorter 
than the characteristic MHD time $\tau_{MHD} \sim L/V_{Ti}$, where 
$V_{Ti} = \sqrt{2 k_B T_i / m_i}$ is the ion thermal velocity.

For Polomac (\cite{Elio2024}):
\begin{itemize}
    \item ECR heating frequency: $f_{ce} = \SI{4}{\giga\hertz}$ 
          $\Rightarrow \omega_{ce} = 2\pi f_{ce} = \SI{2.51e10}{\radian\per\second}$
    \item $V_{Ti} = \sqrt{2 \times \SI{100}{\electronvolt} \times \SI{1.6e-19}{\joule\per\electronvolt} / \SI{1.67e-27}{\kilo\gram}} = \SI{4.4e4}{\meter\per\second}$
    \item $\tau_{MHD} = L/V_{Ti} = \SI{0.30}{\meter} / \SI{4.4e4}{\meter\per\second} = \SI{6.8e-6}{\second}$
    \item $\omega_{MHD}/\omega_{ce} \approx \num{3.7e-5} \ll 1$
\end{itemize}




\item Quasi neutrality 

We assume the plasma to be quasi-neutral, i.e. that the typical plasma length scale $L$ is much greater than the Debye screening distance $\lambda_D$, and outside of that screening distance $\frac{\rho}{ne} = \frac{n_e - n_i}{n_e} \ll 1$, i.e. there aremany particles in the Debye sphere.

For Polomac (\cite{Elio2024}):
\begin{itemize}
    \item Hydrogen density: $n = \SI{e20}{\per\meter\cubed}$
    \item Ion temperature: $T_i = \SI{100}{\electronvolt}$ (assuming $T_e \approx T_i$)
    \item Debye length: $\lambda_D = \SI{7430}{\meter} \sqrt{\frac{T_{\text{eV}}}{n}} = \SI{7430}{\meter} \sqrt{\frac{100}{\num{e20}}} = \SI{7.4e-6}{\meter}$
    \item $\frac{L}{\lambda_D} = \frac{\SI{1}{\meter}}{\SI{7.4e-6}{\meter}} \approx \num{1.4e5} \gg 1$
    \item Number of particles in Debye sphere: $N_D = n \cdot \frac{4\pi}{3} \lambda_D^3 = \num{e20} \cdot \frac{4\pi}{3} \cdot (\SI{7.4e-6}{\meter})^3 \approx \num{1.7e4} \gg 1$ 
\end{itemize}

This is again connected to the fact that we assume electrons can move very fast.

\item Ideal MHD - zero resistivity ($\eta = 0$)

The classical Spitzer resistivity \cite{Chen2016} is given by:
\begin{equation*}
    \eta_\parallel = 5.2 \times 10^{-5} \frac{Z \ln\Lambda}{T_{eV}^{3/2}} \quad [\Omega\text{-m}]
\end{equation*}

For Polomac (\cite{Elio2024}):
\begin{itemize}
    \item Hydrogen plasma: $Z = 1$
    \item Electron temperature: $T_e = \SI{100}{\electronvolt}$ (assuming $T_e \approx T_i$)
    \item Electron density: $n_e = \SI{e20}{\per\meter\cubed}$
    \item Coulomb logarithm: $\ln\Lambda = 14.3$
    \item Parallel resistivity: 
    \begin{equation*}
        \eta_\parallel = 5.2 \times 10^{-5} \times \frac{1 \times 14.3}{(100)^{3/2}} \approx \SI{7.4e-7}{\ohm\meter}
    \end{equation*}
    \item For comparison: this is comparable to stainless steel ($\eta \approx \SI{7e-7}{\ohm\meter}$) and much more resistive than copper ($\eta \approx \SI{2e-8}{\ohm\meter}$).
\end{itemize}

Since $\eta_\perp \sim 2 \eta_\parallel$ (\cite{Chen2016}) we have $\eta \sim \SI{e-6}{\ohm\meter}$.

{\color{magenta}
$\eta\to 0$ is valid if the Lundquist number is high.
}

\item Neglect E.M. waves:

Current density $\vec{J}$ dominates over displacement current $\epsilon_0 \frac{\partial \vec{E}}{\partial t}$.
This holds when plasma timescales are much longer than electromagnetic wave periods, and characteristic MHD wave velocities are much slower than the speed of light.

Assuming low resisitivity (so $E \sim v \wedge B$), $\partial_t \sim \frac{1}{\tau} \sim \frac{v}{L}$, $\nabla \wedge B \sim \frac{B}{L}$ and using $c^2 = \frac{1}{\epsilon_0 \mu_0}$
we can compare the displacement current to the curl of the magnetic field and see that its contribution is small.

The characteristic wave speed in MHD is the Alfvén velocity (\cite{Freidberg2014, Chen2016}):
\[
v_A = \frac{B}{\sqrt{\mu_0 \varrho}}
\]

We get for Polomac (\cite{Elio2024}):
\begin{itemize}
    \item Mass density: $ \varrho = n m_i = \SI{e20}{\per\meter\cubed} \times \SI{1.67e-27}{\kilo\gram} = \SI{1.67e-7}{\kilo\gram\per\meter\cubed}$
    \item Magnetic field: $B = \SI{0.25}{\tesla}$ (midpoint of \SIrange{0.2}{0.3}{\tesla} operating range)
    \item Alfvén velocity: $v_A = \frac{\SI{0.25}{\tesla}}{\sqrt{4\pi \times \SI{e-7}{\henry\per\meter} \times \SI{1.67e-7}{\kilo\gram\per\meter\cubed}}} = \SI{5.5e4}{\meter\per\second}$
    \item Speed of light: $c = \SI{3e8}{\meter\per\second}$
    \item $\frac{v_A}{c} = \frac{\SI{5.5e4}{\meter\per\second}}{\SI{3e8}{\meter\per\second}} \approx \num{1.8e-4} \ll 1$
\end{itemize}

Since MHD wave velocities are non-relativistic ($v_A \ll c$), we can safely neglect the displacement current $\epsilon_0 \partial \vec{E}/\partial t$ compared to $\vec{J}$ in Ampère's law.

Note that $\mu_0 \vec{J} = \nabla \wedge \vec{B}$ implies $\nabla \cdot \vec{J} = 0$.


\item Small Larmor radius $\left(\frac{L}{\rho_i} \gg 1\right)$

The particles are tightly bound to the electric field lines. Inserting the definition of the larmor radius this leads to
\[
\left|\vec{V} \wedge \vec{B}\right| \gg \left|\frac{1}{en} \frac{\partial}{\partial \vec{r}} \cdot \overleftrightarrow{P_e}\right|
\]

Equivalently without the larmor radius
\[
\frac{\vec{V}\wedge\vec{B}}{\frac{1}{en} \frac{\partial}{\partial \vec{r}} \cdot \overleftrightarrow{P_e}}
\sim \frac{v_{th, i} B}{\frac{n_e k_B T_e}{e n L}}
\]

We see that this is negligible for the polomac's values (\cite{Elio2024}).



\item Isotropic pressure and incompressible motion

    \textbf{Isotropic pressure:} The ideal MHD model assumes that the pressure tensor is isotropic, 
    $\overleftrightarrow{P}^{CM} = P$, where $P$ is a scalar pressure. This assumption 
    is not rigorously justified for low-collisionality fusion plasmas, where pressure can become 
    anisotropic ($P_\parallel \neq P_\perp$). However, as shown in \cite{Freidberg2014} Section 2.4.3, 
    MHD stability predictions remain qualitatively accurate because the most unstable modes couple 
    to incompressible motions where the shortcomings of the isotropic pressure assumption have 
    minimal impact.

    \textbf{Adiabatic closure:}
    From the ideal gas law, we see that the adiabatic equation of state ($\frac{ds}{dt}=0$) is equivalent
     to $\frac{d}{dt}(\frac{P}{\varrho^\gamma})$ (see \autoref{chap: App. Adiabatic Equation of State}).
    
    and mass conservation, we obtain
    \[
        \frac{d}{dt}\left(\frac{P}{\varrho^\gamma}\right) = 0 \iff \frac{dP}{dt} + \gamma P \nabla \cdot \vec{V} = 0
    \]
    
    Because particles can move freely along $\vec{B}$ but are bound to field lines by gyromotion, 
    a more physically accurate model would assume rapid heat flow along magnetic field lines 
    ($\vec{B}\cdot\nabla T = 0$) rather than adiabatic behavior.

    \textbf{Incompressible approximation:} For incompressible motions ($\nabla \cdot \vec{V} = 0$) 
    and perfect conductivity ($\eta = 0$), both the adiabatic closure and the infinite parallel 
    thermal conductivity model reduce to the same result (see \cite{Freidberg2014}, Section 2.4.4):
    \[
    \frac{dP}{dt} = 0
    \]

    An analysis in Chapter 8.9 of \cite{Freidberg2014} shows that geometries with closed field lines 
    (such as the Polomac) support both compressible and incompressible instabilities, but the most 
    unstable modes are typically incompressible. For symmetry-preserving modes that cannot be fully 
    incompressible, quantitative stability predictions using ideal MHD should be viewed with caution, 
    though qualitative trends remain valid.
    Because the isothermal equation of state (from ideal gas law) is $\frac{d}{dt}(\frac{P}{\varrho})$,
    so similar but with $\gamma=1$, we conclude that whenever a result depends on $\gamma$, isothermal
    and adiabatic assumptions produce different outputs and we have to be careful with the result.


    In static equilibrium ($\vec{V} = 0$), the incompressibility condition is automatically satisfied.

    If more accurate models are required for the fusion-relevant regime, kinetic MHD or double 
    adiabatic (CGL) theory can be used. Details can be found in \cite{Freidberg2014}.





\end{enumerate}



\subsection{Simplifications}


\begin{itemize}


\item Charge continuity equation:

Charge continuity reduces to $\nabla \cdot \vec{J} = 0$, which is automatically 
satisfied by Ampere's law (Ass. 4: $\nabla \cdot (\nabla \wedge \vec{B}) = 0$). It is thus obsolete.

\[
\underbrace{\frac{\partial \rho}{\partial t}}_{\text{Ass. 2}} + \underbrace{\frac{\partial}{\partial \vec{r}} \cdot \vec{J} = 0}_{\text{Ass. 4}}
\]

\item Momentum equation:
\[
\varrho \frac{d\vec{V}}{dt} = \underbrace{\rho \vec{E}}_{\text{Ass. 2}} + \vec{J} \wedge \vec{B} - \frac{\partial}{\partial \vec{r}} \cdot \underbrace{\overleftrightarrow{P}^{CM}}_{\text{Ass. 6}}
\]

\item Ohm's law:
\[
\underbrace{\frac{m_e}{n e^2} \left[\frac{\partial \vec{J}}{\partial t} + \frac{\partial}{\partial \vec{r}} \cdot (\vec{V}\vec{J} + \vec{J}\vec{V})\right]}_{\text{Ass. 1}} = \vec{E} + \vec{v} \wedge \vec{B} - \underbrace{\frac{1}{en} \frac{\partial}{\partial \vec{r}} \cdot \overleftrightarrow{P^{CM}_e}}_{\text{Ass. 5}}- \underbrace{\eta \vec{J}}_{\text{Ass. 3}}
\]
\end{itemize}



\begin{itemize}

\item Maxwell equations
\begin{align*}
\underbrace{\nabla \cdot \vec{E} = \frac{\rho}{\epsilon_0}}_{\text{redundant}} & & \nabla \wedge \vec{E} &= -\frac{\partial \vec{B}}{\partial t} \\
\nabla \cdot \vec{B} &= 0 & \nabla \wedge \vec{B} &= \mu_0 \left(\vec{J} + \underbrace{\epsilon_0 \frac{\partial \vec{E}}{\partial t}}_{\text{Ass. 4}}\right)
\end{align*}
\end{itemize}





\subsection{The ideal MHD equations}
The ideal magnetohydrodynamic equations consist of
\begin{itemize}
    \item the mass-continuity equation:
            \begin{equation}
                \frac{\partial\varrho}{\partial t} + \nabla(\varrho \vec{V}) = 0
            \end{equation}
    \item the momentum equation:
            \begin{equation}
                \varrho \frac{d \vec{V}}{dt} = \vec{J} \wedge \vec{B} - \nabla P
            \end{equation}
    \item Ohm's law:
            \begin{equation}
                \vec{E} + \vec{V} \wedge \vec{B} = 0
            \end{equation}
    \item Closure equation:
            \begin{equation}
                \frac{d}{dt}\left(P \varrho^{-\gamma}\right) = 0
            \end{equation}
    \item Maxwell equations:
            \begin{equation}
                \nabla\wedge\vec{E} = - \frac{\partial \vec{B}}{\partial t}
            \end{equation}
            \begin{equation}
                \nabla\wedge\vec{B} =  \mu_0 \vec{J}
            \end{equation}
            \begin{equation}
                \nabla\cdot\vec{B} = 0
            \end{equation}
\end{itemize}